\documentclass[12pt]{article}
\usepackage[utf8]{inputenc}
\usepackage[T1]{fontenc}
\usepackage{hyperref}
\title{Bugan Language Resources for K-12 Curriculum Development}
\author{Compiled by Claude Sonnet 5.5 (claude-sonnet-5-5) for Andr\'as Kornai}
\date{}
\begin{document}
\maketitle

\section{Introduction}
The code bgh (Glottocode buga1251) is listed in the task as ``Bugan (Retired)''. The searches indicate that the language called Bugan (also Bogan, Pakan, Bugeng) is catalogued under ISO 639-3 \emph{bbh}, not bgh; the bgh entry appears to be a retired code, and the family label ``Bookkeeping'' in the task is a catalogue artefact. The resources below concern the Bugan language of Wenshan Prefecture, Yunnan, China (Austroasiatic, probably Pakanic), with roughly 2,700--3,000 speakers, most in southern Guangnan and northern Xichou counties. Many speakers also use Southwestern Mandarin; no established orthography was reported.

\section{Primary Educational Resources}
No primers, textbooks or school programmes for Bugan were found.

\section{Linguistic and Reference Materials}
\begin{itemize}
\item Wikipedia, Bugan language: \url{https://en.wikipedia.org/wiki/Bugan_language}
\item Austronesian Basic Vocabulary Database, Bugan (Nala) word list: \url{https://abvd.eva.mpg.de/austronesian/language.php?id=961}
\item ABVD, Bugan (Manlong) word list: \url{https://lpan.eva.mpg.de/austronesian/language.php?id=962}
\item Omniglot, Bugan: \url{https://mail.omniglot.com/writing/bugan.htm}
\end{itemize}

\section{Religious and Community Resources}
\begin{itemize}
\item Asia Harvest people-group profile (PDF): \url{https://people-groups.asiaharvest.org/people-group-profiles/yunnan/bugan.pdf}
\end{itemize}

\section{Digital and Audio Resources}
\begin{itemize}
\item Endangered Languages Project entry: \url{https://endangeredlanguages.com/lang/1238}
\end{itemize}

\section{Curriculum Development Context}
Bugan has no established orthography and no teaching materials were found. Any curriculum would depend on first developing a writing system with the community; Mandarin is the practical language of schooling.

\end{document}